\documentclass{article}

\usepackage[margin=2.5cm]{geometry}

\usepackage{graphicx}
\usepackage{float}
\usepackage{booktabs}
\usepackage{listings}
\usepackage{xcolor}
\usepackage{hyperref}
\usepackage{subcaption}
\usepackage{comment}
\usepackage{parskip}
\usepackage{setspace}
\usepackage{amsmath}
\usepackage{amssymb}
\usepackage{array}
\usepackage{tikz}
\usepackage{markdown}
\usepackage{fancyvrb}


\newcounter{plantumldiagram}

\newenvironment{plantuml}[3]
{%
    \stepcounter{plantumldiagram}%
    \def\plantumlwidth{#1}%
    \def\plantumlcaption{#2}%
    \def\plantumllabel{#3}%
    \VerbatimEnvironment
    \begin{VerbatimOut}{PlantUML/generated/\jobname-plantuml-\theplantumldiagram.puml}%
}
{%
    \end{VerbatimOut}%
    \immediate\write18{java -jar "PlantUML/plantuml.jar" -tpng - scale 4 "PlantUML/generated/\jobname-plantuml-\theplantumldiagram.puml"}%
    \begin{figure}[H]
        \centering
        \includegraphics[width=\plantumlwidth]{PlantUML/generated/\jobname-plantuml-\theplantumldiagram.png}
        \caption{\plantumlcaption}
        \label{\plantumllabel}
    \end{figure}%
}

% Change scale to change the fidelity of image

\title{3.2 Models}
\author{Student Number: 4006409}
\date{\today}



\begin{document}

\maketitle
\section{Engineering Models and Mathematical Analysis}

The primary design scenario, S3 — Emergency Vehicle Priority, was evaluated using a combination of traffic-flow, signal-timing, control-loop, emergency-priority, fault-response and reliability models. The Python simulation uses a 1-second simulation step over a 30-minute peak-period scenario. The engineering parameters are based on the assumptions implemented in the simulation.

\subsection{Traffic Flow and Queueing Model}

Each intersection approach is represented by an aggregate vehicle queue. Vehicle arrivals are modelled using a time-varying Poisson process. During a green phase, vehicles are discharged at the assumed saturation flow rate.

The queue update is represented by:

$$
Q(t+\Delta t)=Q(t)+A(t)-D(t)
$$

where:
* \(Q(t)\) = queue length at time \(t\)
* \(A(t)\) = number of arriving vehicles during the time step
* \(D(t)\) = number of departing vehicles
* \(\Delta t=1\) s.

The assumed saturation flow is:

$$
s=1900\text{ veh/h/lane}
$$

Converting this to vehicles per second:

$$
s=\frac{1900}{3600}
$$

$$
\boxed{s=0.528\text{ veh/s}}
$$

Therefore, if an approach has a queue and remains green for 10 seconds, its theoretical maximum discharge is:

$$
D=s\times t
$$

$$
D=0.528\times10=5.28\text{ vehicles}
$$

Thus, approximately 5.28 vehicles can be discharged per lane during a 10-second green interval, subject to the number of vehicles actually waiting.

The simulation uses an average demand of approximately 0.21 veh/s per approach. The corresponding flow ratio is:

$$
y=\frac{q}{s}
$$

$$
y=\frac{0.21}{0.528}
$$

$$
\boxed{y\approx0.398}
$$

For the two opposing phase pairs, the combined critical flow ratio used by the model is approximately:

$$
Y=0.398+0.398
$$

$$
\boxed{Y\approx0.796}
$$

This remains below the theoretical saturation condition \(Y=1\), allowing the fixed-time and adaptive control strategies to be compared without demand continually exceeding the available discharge capacity.

\subsection{Fixed-Time Signal Timing — Webster Calculation}

The fixed-time controller provides the baseline against which the adaptive controller is compared. The cycle length is calculated using Webster's optimal-cycle formula:

$$
C_o=\frac{1.5L+5}{1-Y}
$$

There are two signal phases and 4 seconds of lost time per phase:

$$
L=2\times4
$$

$$
\boxed{L=8\text{ s}}
$$

Using \(Y\approx0.796\):

$$
C_o=\frac{1.5(8)+5}{1-0.796}
$$

$$
C_o=\frac{17}{0.204}
$$

$$
\boxed{C_o\approx83.3\text{ s}}
$$

The effective green time available after lost time is:

$$
G=C_o-L
$$

$$
G=83.3-8
$$

$$
\boxed{G\approx75.3\text{ s}}
$$

For approximately equal critical flow ratios, the green time is divided approximately equally between the two phases:

$$
g_i=G\frac{y_i}{Y}
$$

$$
g_i=75.3\frac{0.398}{0.796}
$$

$$
\boxed{g_i\approx37.7\text{ s}}
$$

Therefore, the fixed-time baseline uses an approximately 83-second cycle with approximately 38 seconds of effective green allocated to each phase.

This provides a consistent reference strategy because the timing is calculated from the average demand and is not subsequently modified in response to short-term traffic changes.

\subsection{Adaptive Vehicle-Actuated Control}

The adaptive controller uses real-time traffic information rather than continuously recalculating a complete Webster cycle. The controller applies three main decisions:

1. **Extend** — continue the green phase when vehicles remain queued.
2. **Gap-out** — terminate the phase when the queue clears.
3. **Max-out** — terminate the phase when the maximum green time is reached.

The design parameters are:

$$
G_{min}=10\text{ s}
$$

$$
G_{max}=60\text{ s}
$$

Therefore:

$$
10\leq G_i\leq60\text{ s}
$$

If an approach has no vehicles waiting, its phase can be skipped. This allows unused green time to be transferred to approaches experiencing demand.

The purpose of this control strategy is therefore not simply to reduce the nominal cycle length. Instead, it dynamically allocates green time according to the observed traffic demand.

\subsection{Sensing and Control-Loop Latency}

The closed-loop controller consists of sensing, condition detection and transmission of revised timing information.

The simulation models the following:

* Data update interval: \(U\sim Uniform(3,5)\) s
* Detection latency: \(|N(1.0,0.4)|\) s, capped at 2.5 s
* Transmission latency: \(|N(0.9,0.4)|\) s, capped at 2.5 s.

For the data-update interval:

$$
3\leq U\leq5\text{ s}
$$

Therefore the maximum permitted update interval is:

$$
U_{max}=5.0\text{ s}
$$

which directly corresponds to SR-F01.

For the detection process, the nominal mean is approximately:

$$
\mu_{detect}=1.0\text{ s}
$$

against the requirement:

$$
T_{detect}\leq2.0\text{ s}
$$

The latest simulation produced a mean detection time of approximately 0.99 s and a worst-case value of 2.47 s. Consequently, the mean satisfies the requirement, while the observed worst case exceeds the 2-second limit. This identifies SR-F02 as a performance margin issue requiring further investigation rather than treating the mean value alone as sufficient evidence.

\subsection{Emergency Vehicle Priority}

Emergency priority follows the design precedence:

$$
\boxed{\text{Safety}>\text{Compliance}>\text{Emergency}>\text{Efficiency}}
$$

This is particularly important for conflict C1, where an emergency request occurs while pedestrian clearance is active.

The pedestrian clearance interval is:

$$
T_{ped}=7\text{ s}
$$

If an emergency request occurs during this interval, the request is deferred until the clearance period has completed.

For example, if an emergency request occurs with \(t_r\) seconds remaining in the clearance interval, the minimum additional waiting time is:

$$
T_{wait}=t_r
$$

Once:

$$
T_{ped}=0
$$

the emergency priority request can be granted.

This prevents the pedestrian interval from being interrupted. Consequently:

$$
N_{pedestrian\ violations}=0
$$

and the simulation recorded:

$$
\boxed{SR\text{-}N01=0\text{ violations}}
$$

The emergency-request calculation is:

$$
P_{grant}=
\frac{N_{granted}}{N_{requests}}\times100
$$

For the simulation:

$$
P_{grant}=
\frac{N_{granted}}{N_{requests}}\times100
=100\%
$$

Thus:

$$
\boxed{SR\text{-}F04=100\%}
$$

These results are primarily a verification of the implemented safety rule because emergency requests are deliberately deferred rather than rejected.

\subsection{Fault Detection and Fallback}

Fault response is modelled using three independent latency components:

$$
T_{detect}\sim |N(4.0,2.0)|
$$

$$
T_{fallback}\sim |N(4.5,2.2)|
$$

$$
T_{notify}\sim |N(14.0,6.0)|
$$

The values are capped at 9.8 seconds for detection and fallback and 29 seconds for notification so that the simulated system can be tested against its specified limits.

The requirements are:

$$
T_{detect}\leq10\text{ s}
$$

$$
T_{fallback}\leq10\text{ s}
$$

$$
T_{notify}\leq30\text{ s}
$$

The latest simulation produced:

$$
T_{detect,worst}=8.20\text{ s}
$$

$$
T_{fallback,worst}=5.23\text{ s}
$$

$$
T_{notify,worst}=20.16\text{ s}
$$

Therefore:

$$
8.20<10
$$

$$
5.23<10
$$

$$
20.16<30
$$

All three fault-response timing requirements are satisfied in the simulation.

Once fallback is activated, the affected intersection returns to the safe fixed-time strategy. This ensures that a single sensor or communications failure does not result in loss of traffic control.

\subsection{Reliability and Availability Calculation}

System availability is estimated using the standard steady-state relationship:

$$
A=\frac{MTBF}{MTBF+MTTR}
$$

The current simulation assumes:

$$
MTBF=1500\text{ h}
$$

$$
MTTR=3\text{ h}
$$

Therefore:

$$
A=\frac{1500}{1500+3}
$$

$$
A=\frac{1500}{1503}
$$

$$
A=0.998004
$$

Converting to a percentage:

$$
A=0.998004\times100
$$

$$
\boxed{A\approx99.80\%}
$$

The availability requirement is:

$$
A\geq99\%
$$

and therefore:

$$
99.80\%>99.00\%
$$

so the reliability requirement is satisfied by the analytical model.

The relationship can also be rearranged to estimate MTBF from an assumed availability:

$$
MTBF=\frac{A\times MTTR}{1-A}
$$

For example, using a Victorian traffic-signal availability benchmark of approximately 99.8% and the assumed 3-hour restoration time:

$$
MTBF=\frac{0.998\times3}{1-0.998}
$$

$$
MTBF=\frac{2.994}{0.002}
$$

$$
\boxed{MTBF\approx1497\text{ h}}
$$

which is approximately 1500 hours and provides the basis for the simulation assumption.

Importantly, the 1500-hour MTBF and 3-hour MTTR should be described as **engineering assumptions/calibration values**, rather than claiming they are measured average values for Victorian roadside controllers. The availability calculation is analytical and is separate from the discrete-event traffic simulation.

\subsection{Queue-Reduction Performance Calculation}

The adaptive and fixed-time controllers are compared using their average simulated queue lengths.

The latest simulation produced:

$$
\bar Q_{fixed}=17.06\text{ vehicles}
$$

$$
\bar Q_{adaptive}=13.85\text{ vehicles}
$$

The percentage reduction is calculated as:

$$
\text{Reduction}=
\frac{\bar Q_{fixed}-\bar Q_{adaptive}}
{\bar Q_{fixed}}\times100
$$

Substituting the simulation results:

$$
\text{Reduction}=
\frac{17.06-13.85}{17.06}\times100
$$

$$
=\frac{3.21}{17.06}\times100
$$

$$
\boxed{\text{Reduction}\approx18.82\%}
$$

Thus, under the simulated traffic demand, the adaptive controller reduced the average queue length by approximately **18.8% relative to the fixed-time baseline**.

This result represents the primary efficiency outcome of the simulation. It demonstrates that responding to current queue conditions can reduce average queued vehicles compared with a fixed-time plan calculated from average demand.


\end{document}